电子初态的有限纵向带宽则通过输入密度矩阵进入。令
\begin{equation*}
 \rho_\ee^{\rm in}
 =\int\!\dd p\,\dd p'\,
 \rho_e(p,p')|p\rangle\langle p'|,
\end{equation*}
并用总散射算符 $S$ 定义电子约化输出
\begin{equation}
{
 \rho_\ee^{\rm out}
 =\operatorname{Tr}_{\rm ph}
 \left[S(\rho_\ee^{\rm in}\otimes\rho_{\rm ph})S^\dagger\right].}
\label{eq:finite_envelope_reduced_output_dp47}
\end{equation}
取能量归一化的电子基 $\langle E|E'\rangle=\delta(E-E')$；若还存在未分辨自旋或横向通道，下式的 $p_e(E)$ 对这些指标求和。电子的理想能谱密度为 $p_e(E)=\langle E|\rho_\ee^{\rm out}|E\rangle$，单位为能量的倒数，满足 $\int p_e(E)\dd E=1$。令 $R_E(\varepsilon-E)\ge0$ 表示真实能量 $E$ 被读成 $\varepsilon$ 的条件密度，$\int R_E(x)\dd x=1$。谱仪记录的密度先由卷积给出，再在互不相交的像素能窗 $I_j$ 内积分为无量纲概率：
\begin{equation}
{
 \begin{aligned}
 p_{\rm det}(\varepsilon)&=\int\dd E\,R_E(\varepsilon-E)p_e(E),\\
 P_j^{\rm meas}&=\int_{I_j}\dd\varepsilon\,p_{\rm det}(\varepsilon).
 \end{aligned}}
\label{eq:pinem_detector_convolution_dp47}
\end{equation}
卷积后的 $p_{\rm det}(E_j)$ 仍是密度，不能直接称为第 $j$ 个像素的概率；窄像素近似为 $P_j^{\rm meas}\simeq p_{\rm det}(E_j)\Delta E_j$。当像素窗覆盖全部读出能量时，交换两次非负积分便得 $\sum_jP_j^{\rm meas}=1$；有限探测范围要另保留漏出概率。后续统计反演应使用\eqnrefs{eq:pinem_detector_convolution_dp47}的像素核，而不是直接把理想 $P_\ell$ 与实验直方图比较。

有限输入能散 $\sigma_E$ 是否会洗掉相邻边带取决于两个不同条件。分辨边带需要
\begin{equation}
 \epsilon_{\rm res}\equiv\frac{\sigma_E}{\hbar\omega}\ll1,
\end{equation}
而保持输入相干的相位条件由输入密度矩阵非对角宽度 $\sigma_{E,{\rm coh}}$ 决定。二者一般不相等；能谱窄并不自动意味着长相干时间。
